% Part: second-order-logic
% Chapter: syntax-and-semantics
% Section: introduction

\documentclass[../../../include/open-logic-section]{subfiles}

\begin{document}

\olfileid{sol}{syn}{int}

\olsection{تعارف}

پہلے درجے دی منطق وچ اسیں کسے زبان دیاں غیر منطقی
علامتاں، یعنی اوہدے !!{constant}s، !!{function}s تے
!!{predicate}s، نوں منطقی علامتاں نال ملا کے پہلے درجے
دیاں !!{structure}s بارے گل ظاہر کردے آں۔ ایس کم لئی پورا اترن دا تصور ورتیدا اے۔ ایہ !!a{structure}~$\Struct{M}$
تے متغیراں دی قدر بندی~$s$ نوں !!a{formula}~$!A$ نال جوڑدا اے۔
$\Sat{M}{!A}[s]$ اُدوں تے صرف اُدوں لاگو اے جدوں $!A$ دی کہی گل
سچ ہووے، جدوں اوہدے
!!{constant}s، !!{function}s تے !!{predicate}s دی
تعبیر $\Struct{M}$ دے مطابق تے آزاد متغیراں دی
تعبیر~$s$ دے مطابق ہووے۔ !!{identity}~$\eq$ دی
تعبیر $\Sat{M}{!A}[s]$ دی تعریف وچ پہلے توں شامل
اے،جیویں~$\lforall$ تے $\lexists$ دی تعبیر وی۔
پہلی علامت نوں ہمیشہ ساخت دے !!{domain}~$\Domain{M}$
اُتے برابری دے تعلق دے طور تے سمجھیا جاندا اے، تے مسوّراں دا دائرہ
پورا !!{domain} ہوندا اے۔ پر خاص گل ایہ اے کہ مسوّریت صرف
!!{domain} دے عناصر اُتے ہو سکدی اے؛ایس لئی مسوّر
دے بعد صرف شے دے !!{variable}s آ سکدے نیں۔

دوجے درجے دی منطق وچ زبان تے پورا اترن دی تعریف
دونوں نوں وسیع کیتا جاندا اے،تاکہ آزاد تے پابند
فنکشنی تے محمولی متغیر،تے اوہناں اُتے مسوّریت،شامل
ہوون۔ ایہہ متغیر !!{function}s تے !!{predicate}s نال اوہو تعلق رکھدے نیں
جیہڑا شے دے متغیر !!{constant}s نال رکھدے نیں۔ دوجے درجے دی
منطق دے ٹرم تے
!!{formula}s بناؤن وچ ایہہ وی اوہی کردار نبھاؤندے
نیں،تے اوہناں اُتے مسوّریت وی ملدی جلُدی طرح
نبھائی جاندی اے۔ \emph{معیاری} معنیات وچ دوجے درجے
دے مسوّر درست نوع دیاں ساریاں ممکنہ چیزاں اُتے
چلدے نیں۔ فنکشنی متغیراں لئی ایہ سارے $n$-رکنی فنکشن نیں،
جیہڑے $\Domain{M}$ دے ترتیب وار ٹپل نوں~$\Domain{M}$ دا اک
عنصر دیندے نیں۔ محمولی متغیراں لئی ایہ سارے $n$-رکنی تعلق نیں۔ مثال لئی
$\lforall[\Obj{v_0}][(\Obj{P^1_0}(\Obj{v_0}) \lor \lnot
  \Obj{P^1_0}(\Obj{v_0}))]$ پہلے تے دوجے دونوں درجیاں
دی منطق دا فارمولا اے؛پر دوجے درجے وچ اسیں
$\lforall[\Obj{V^1_0}][\lforall[\Obj{v_0}][(\Obj{V^1_0}(\Obj{v_0}) \lor
    \lnot \Obj{V^1_0}(\Obj{v_0}))]]$ تے
$\lexists[\Obj{V^1_0}][\lforall[\Obj{v_0}][(\Obj{V^1_0}(\Obj{v_0}) \lor
    \lnot \Obj{V^1_0}(\Obj{v_0}))]]$ وی ویکھ سکدے آں۔
کیوں جے انہاں وچ کوئی آزاد متغیر نہیں، ایس لئی ایہ دوجے درجے
دی منطق دے !!{sentence}s نیں۔ ایتھے $\Obj{V^1_0}$
دوجے درجے دا $1$-رکنی محمولی متغیر اے۔ $\Obj{V^1_0}$
دیاں جائز تعبیراں اوہو نیں جیہڑیاں اسیں کسے
$1$-رکنی !!{predicate}،مثلاً $\Obj{P^1_0}$،نوں دے
سکدے آں،یعنی
~$\Domain{M}$ دے ذیلی سیٹ۔ ایہناں ذیلی سیٹاں اُتے مسوّریت دا مطلب ایہ اے کہ $\lforall[\Obj{v_0}][(\Obj{V^1_0}(\Obj v_0) \lor \lnot
  \Obj{V^1_0}(v_0))]$، $\Domain{M}$ دا ذیلی سیٹ
$\Obj{V^1_0}$ دی قدر دے طور تے دین دے
ہر طریقے لئی، یا گھٹ توں گھٹ اک طریقے لئی، درست ہووے۔
ہر دِتّی شے کسے سیٹ وچ یا شامل ہوندی اے یا شامل نہیں ہوندی۔
ایس لئی دوویں جملے ہر
!!{structure} وچ سچ نیں۔
\end{document}
